\section*{Chapitre \ref{ch:b3:glaciology} --- Glaciologie}

\begin{solution}{exo:b3:glaciology:1}
$\tau_b = \rho gH\sin\alpha = 917\times9.81\times250\times\sin\ang{3} =
\qty{118}{kPa}$, environ \qty{1.2}{bar}~: l'ordre de grandeur des contraintes
sous les glaciers réels.
\end{solution}

\begin{solution}{exo:b3:glaciology:2}
(a)~$\qty{2.9e15}{m^3}\times917 = \qty{2.66e18}{kg} = \qty{2.66e6}{Gt}$~;
divisé par \qty{360}{Gt/mm}, on obtient \qty{7390}{mm}, soit \qty{7.4}{m}, la valeur annoncée.
(b)~$\qty{3.0e16}{m^3}\times917/360 = \qty{76}{m}$, contre \qty{58}{m}.
La glace du Groenland repose surtout sur un socle au-dessus du niveau de la mer, si bien que toute cette glace
ajouterait de l'eau à l'océan. Une grande partie de celle de l'Antarctique repose sur un socle situé sous
le niveau de la mer~: cette glace déplace déjà de l'eau de mer, et sa fonte ne fait que remplacer
l'eau qu'elle déplace~; les plates-formes de glace flottantes n'ajoutent rien du tout.
\end{solution}

\begin{solution}{exo:b3:glaciology:3}
Gain $\num{6e6}\times1.2 = \qty{7.2e6}{m^3}$, perte $\num{4e6}\times2.1 =
\qty{8.4e6}{m^3}$ d'eau~: bilan net \qty{-1.2e6}{m^3}, soit un bilan spécifique moyen
de $-1.2\times10^6/10^7 = \qty{-0.12}{m}$ d'équivalent en eau. Le
rapport de la surface d'accumulation est $6/10 = 0.6$. Le glacier perd de la masse~:
il se réduit.
\end{solution}

\begin{solution}{exo:b3:glaciology:4}
$1 - 820/917 = \num{0.106}$ et $1 - 830/917 = \num{0.095}$~: environ un dixième
du volume reste occupé par des pores lorsque ceux-ci se referment.
\end{solution}

\begin{solution}{exo:b3:glaciology:5}
$\tau_b = 917\times9.81\times300\times\sin\ang{3} = \qty{141.2}{kPa}$. D'après
le théorème de la dalle de glace, $u_s = \tfrac{2A}{4}\tau_b^3H = 0.5\times\num{1.43e-16}
\times(\num{1.412e5})^3\times300 = \qty{60}{m/an}$. Épaisseur doublée~: $u_s
\propto H^{n+1} = H^4$, seize fois plus rapide. Pente divisée par deux~: $u_s \propto
(\sin\alpha)^3$, huit fois plus lente.
\end{solution}

\begin{solution}{exo:b3:glaciology:6}
À \qty{-25}{\celsius} (\qty{248.15}{K}), $A = \num{3.61e-13}\exp(-60000/
(8.314\times248.15)) = \qty{8.5e-26}{Pa^{-3}.s^{-1}}$, soit
\qty{2.7e-18}{Pa^{-3}.an^{-1}}. À \qty{-10}{\celsius}, $\num{3.61e-13}
\exp(-27.42) = \qty{4.4e-25}{Pa^{-3}.s^{-1}} = \qty{1.4e-17}{Pa^{-3}.an^{-1}}$.
À \qty{0}{\celsius} (branche tempérée) \qty{1.4e-16}{Pa^{-3}.an^{-1}}. Sous
les mêmes contraintes, les vitesses sont dans le rapport des \omterm{def:b3:glaciology:glen}{coefficients de fluage}~:
$\num{1.43e-16}/\num{2.67e-18} = 54$.
\end{solution}

\begin{solution}{exo:b3:glaciology:7}
$h = \sqrt{2\times10^5\times10^6/(917\times9.81)} = \qty{4.7}{km}$, proche de
\qty{4.9}{km} de l'Antarctique~: le modèle plastique donne l'épaisseur
d'une calotte glaciaire à partir de sa seule largeur.
\end{solution}

\begin{solution}{exo:b3:glaciology:8}
$\tau = 200/5 = \num{40}\text{ ans}$. La longueur continuerait de s'ajuster pendant
plusieurs décennies, soit quelques temps de réponse, après l'arrêt du réchauffement~:
les glaciers actuels n'ont pas encore rattrapé le climat d'aujourd'hui.
\end{solution}

\begin{solution}{exo:b3:glaciology:9}
(a)~$h = \sqrt{2\times2.03\times20\times\num{8.64e6}/(917\times\num{3.34e5})}
= \qty{1.51}{m}$. Puisque $h^2 \propto t$, une épaisseur de \qty{3}{m} demande $100\times
(3/1.51)^2 = 393$ jours, soit plus qu'un hiver. (b)~La glace bien plus épaisse que
\qty{3}{m} ne peut pas se former par gel en un seul hiver~: c'est de la glace qui a
survécu à plusieurs hivers, ou de la glace entassée par le vent et les courants dans des crêtes.
\end{solution}

\begin{solution}{exo:b3:glaciology:10}
(a)~$H = (1025/917)\times800 = \qty{894}{m}$. (b)~À \qty{900}{m}, l'épaisseur
à la \omterm{def:b3:glaciology:grounding-line}{ligne d'échouage} vaut \qty{1006}{m}, soit $9/8$ de la valeur précédente~; le flux
augmente de $(9/8)^5 = 1.8$. (c)~Lorsque la ligne recule sur un socle qui s'enfonce
vers l'intérieur des terres, la glace au niveau de la ligne est plus épaisse, débite davantage, amincit la glace
en amont, et la ligne recule davantage~: une rétroaction positive. Une plate-forme de glace
qui retient la glace posée sur le socle réduit le flux à la \omterm{def:b3:glaciology:grounding-line}{ligne d'échouage} et peut
la maintenir en place même sur une pente descendante.
\end{solution}

\begin{solution}{exo:b3:glaciology:11}
Avec $h(r) = c\sqrt{R - r}$, $c = \sqrt{2\tau_0/(\rho g)}$, $V =
\int_0^R2\pi rc\sqrt{R - r}\,\mathrm dr$. En posant $s = R - r$,
\[
\int_0^R(R - s)\sqrt s\,\mathrm ds = \tfrac23R^{5/2} - \tfrac25R^{5/2} =
\tfrac4{15}R^{5/2},
\]
d'où $V = \frac{8\pi}{15}cR^{5/2}$. Doubler $R$
multiplie $V$, ainsi que l'\omterm{def:b3:glaciology:sle}{équivalent en niveau marin}, par $2^{5/2} = 5.7$.
\end{solution}

\begin{solution}{exo:b3:glaciology:12}
(a)~Mars~: $14.13/16.34 - 1 = \qty{-13.5}{\percent}$~; septembre~: $4.81/7.05
- 1 = \qty{-31.8}{\percent}$. (b)~En mars, le bord de la banquise est retenu essentiellement par
les côtes qui entourent l'océan Arctique, et la recongélation chaque hiver
remplit de nouveau le bassin~; la glace plus mince des dernières années fond plus
complètement chaque été. La rétroaction glace--albédo agit en été~: l'eau libre
absorbe environ \qty{94}{\percent} du rayonnement solaire, la glace en réfléchit
\qtyrange{50}{70}{\percent}, si bien que chaque perte de glace estivale réchauffe l'eau et
la fait fondre davantage.
\end{solution}

\begin{solution}{pb:b3:glaciology:1}
\textbf{1.} La colonne au-dessus du socle, de poids $\rho gH$ par unité de surface, est
poussée vers l'aval par la composante $\rho gH\sin\alpha$, équilibrée par la
contrainte de cisaillement à sa base~; pour de faibles pentes $\sin\alpha \approx \alpha$ (le
théorème de la dalle de glace avec la pente de surface).

\textbf{2.} $917\times9.81\times580\times0.0190 = \qty{99.1}{kPa}$~;
$917\times9.81\times470\times0.0240 = \qty{101.5}{kPa}$~; $917\times9.81
\times335\times0.0330 = \qty{99.4}{kPa}$.

\textbf{3.} La contrainte reste la même à \qty{2}{\percent} près pendant que
l'épaisseur est presque divisée par deux~: la calotte est proche d'un comportement plastique, avec $\tau_0 \approx
\qty{100}{kPa}$.

\textbf{4.} $2\times10^5/(917\times9.81) = \qty{22.23}{m}$.

\textbf{5.} $h = \sqrt{22.23\,(R - r)}$~: \qty{577}{m}, \qty{472}{m} et
\qty{333}{m} à 15, 10 et \qty{5}{km} de la marge, à quelques mètres
près des sondages.

\textbf{6.} $\sqrt{22.23\times20000} = \qty{667}{m}$.

\textbf{7.} Comme dans l'\cref{exo:b3:glaciology:11}~: $V = \int_0^R2\pi r
\sqrt{22.23(R - r)}\,\mathrm dr = \frac{8\pi}{15}\sqrt{2\tau_0/(\rho
g)}\,R^{5/2}$.

\textbf{8.} $\frac{8\pi}{15}\times4.715\times(\num{2e4})^{5/2} =
\qty{4.47e11}{m^3} = \qty{447}{km^3}$.

\textbf{9.} $\pi\times20^2 = \qty{1257}{km^2}$~; épaisseur moyenne
$447/1257 = \qty{0.356}{km}$, \qty{356}{m}, soit un peu plus de la moitié de
l'épaisseur centrale.

\textbf{10.} $2^{5/2} = 5.7$ fois.

\textbf{11.} $A = \num{3.61e-13}\exp(-60000/(8.314\times263.15)) =
\qty{4.45e-25}{Pa^{-3}.s^{-1}}$, multiplié par \qty{3.156e7}{s}~:
\qty{1.40e-17}{Pa^{-3}.an^{-1}}.

\textbf{12.} À \qty{10}{km}, $H = \qty{471.5}{m}$, pente $471.5/20000 =
\num{0.0236}$, $\tau_b = \qty{100}{kPa}$; $u_s = \tfrac{2A}{4}\tau_b^3H =
0.5\times\num{1.40e-17}\times10^{15}\times471.5 = \qty{3.3}{m/an}$.

\textbf{13.} $q = \tfrac{2A}{5}\tau_b^3H^2 = 0.4\times\num{1.40e-17}
\times10^{15}\times471.5^2 = \qty{1250}{m^2/an}$~; vitesse moyenne $q/H =
\qty{2.6}{m/an}$, soit les quatre cinquièmes de la vitesse de surface.

\textbf{14.} $2\pi rq = \pi r^2\bar b$, donc $\bar b = 2q/r = 2\times1250/10^4
= \qty{0.25}{m}$ de glace par an sur les \qty{10}{km} intérieurs.

\textbf{15.} $h = 450$ là où $R - r = 450^2/22.23 = \qty{9.11}{km}$, soit à
$r = \qty{10.9}{km}$~; la \omterm{def:b3:glaciology:mass-balance}{zone d'accumulation} est le disque qu'il délimite, $(10.9/
20)^2 = 0.30$ de la calotte.

\textbf{16.} $A$ augmente de \num{10.2}, et il en va de même pour les vitesses sous les mêmes
contraintes. La glace s'évacuerait plus vite que ne l'alimente le bilan~: la
calotte s'amincirait et s'étalerait jusqu'à ce que les contraintes, plus faibles, égalent à nouveau le
bilan~; une calotte plastique plus molle a une contrainte seuil plus faible et un
profil plus plat.

\textbf{17.} $\num{4.47e11}\times917 = \qty{4.10e14}{kg} = \qty{410}{Gt}$.

\textbf{18.} $410/360 = \qty{1.14}{mm}$.

\textbf{19.} $356/0.5 = 711$ ans.

\textbf{20.} À mesure que la surface s'abaisse, elle se retrouve dans un air plus chaud et fond plus vite,
la rétroaction altitude--bilan de masse~; et à mesure que la calotte s'amincit, les contraintes et
le régime d'écoulement changent aussi. L'amincissement s'accélère.

\textbf{21.} $410/243 = \num{1.7}\text{ ans}$~: la calotte entière représente moins de deux
ans de perte du Groenland en 2010--2019.

\textbf{22.} La calotte glaciaire d'Arkel représente un \omterm{def:b3:glaciology:sle}{équivalent en niveau marin} de
\textbf{\qty[text-series-to-math=true]{1.1}{mm}} (\qty{1.14}{mm}).
\end{solution}
